% UPC starter; recheck the current contest rules before submission.
% Replace all bracketed fields and inspect the rendered PDF and metadata.
\documentclass[11pt,a4paper]{article}
\usepackage[margin=25mm]{geometry}
\usepackage{amsmath,amssymb,graphicx,booktabs}
\usepackage{fancyhdr}
\usepackage[hidelinks]{hyperref}
\newcommand{\TeamNumber}{[TEAM NUMBER]}
\newcommand{\ProblemLetter}{[A OR B]}
\newcommand{\PaperTitle}{[RESEARCH TITLE]}
\pagestyle{fancy}
\fancyhf{}
\fancyhead[L]{Team \TeamNumber}
\fancyfoot[C]{\thepage}
\setlength{\headheight}{14pt}
\fancypagestyle{summary}{\fancyhf{}\fancyhead[L]{Team \TeamNumber}}
\hypersetup{pdfauthor={},pdfsubject={},pdfkeywords={}}
\begin{document}
\begin{titlepage}
\thispagestyle{summary}
\begin{center}
Problem \ProblemLetter

\vspace{1em}
{\Large\bfseries \PaperTitle\par}
\end{center}
\begin{abstract}
[Replace with a self-contained summary of at most the current word limit
(300 words in the rules verified on 2026-11-07). State the problem,
method and key assumptions, direct answers, principal numerical results
with units, and material conditions. Count the final abstract separately.]
\end{abstract}
\end{titlepage}
\setcounter{page}{2}
\section{Problem and scope}
[Restate the questions, outputs and definitions that affect the answers.]
\subsection{Research contributions}
[State the physical result, prediction method and validation contribution.
Pair each claim with a computed or measured result, a comparison baseline,
and the conditions under which the claim holds.]
\section{Assumptions and physical model}
[Define variables, units, parameter sources, validity and governing relations.]
\section{Methods}
[Explain the actual algorithms, inputs, boundary conditions and tolerances
in the main text; an appendix of code does not replace this explanation.]
\section{Results and validation}
[Give actual answers, baseline comparisons, physical explanations,
appropriate independent checks, and uncertainty with its meaning.]
\subsection{Sensitivity to physical and observational factors}
[Compare changes to timing, position error, force assumptions and model
weights under a fixed prediction target. Report the parameter change,
point displacement, joint region area, and sampling or convergence checks.]
\subsection{Value of an additional observation}
[Keep the original decision time for the main answer. In the extension,
give each input mask and deadline, compare equal computational budgets on
the same shots, and report remaining flight time, error, coverage, and area.]
\section{Discussion and limitations}
[Explain dominant mechanisms, robustness, applicability and failure cases.]
[Distinguish the official example, held-out synthetic results, and field
cases. State which contact timestamps, coordinate calibration, and event
labels are needed to test the method in the field.]
\section{Conclusion}
[Answer each problem requirement within the established limits.]
\begin{thebibliography}{99}
\bibitem{replace} [Insert verified sources and cite them at relevant claims.]
\end{thebibliography}
% Add long derivations/code appendices here, before the AI report.
\clearpage
\section*{Report on the Use of AI}
[When required by the current rules, include the actual system and use,
verbatim prompts and complete visible outputs for relevant interactions.
Convert the report generated by \texttt{ai\_report.py} to the manuscript format and
check the rendered result. This marker is not a disclosure report.
Only if no generative AI was actually used, state that truthfully.]
\end{document}
